% Einheit 3 von 4 – Normalform und p-q-Formel (bank/quadratische-gleichungen/e3.jsonl, zusammenbau v0.1)
%% TODO Zweigzeile Teil 1: was hier gelernt wird (Satz in Schülersprache aus dem Einheitstitel) – Einheit 3
\einheitenkopf[e3]{Einheit 3 von 4 · Normalform und p-q-Formel}
\zweigzeile{ab Kl. 9, je nach Buch bis Kl. 10 · P10 oft}
\setcounter{aufgabe}{15}

%% TODO Titel: Ich-kann-Satz fehlt in der Bank (Platzhalter: Kettenname) – Nr. 16
%% TODO Anweisung: ein Satz über der ersten Teilaufgabe fehlt in der Bank – Nr. 16
\begin{aufgabe}{Welche Form?}
\begin{teile}
\teil $x^2 - 3x + 1 = 0$ – welche Form, welcher Weg?\\ \kreuz{reinquadratisch – Wurzelziehen} \\ \kreuz{Produkt gleich null – Nullprodukt} \\ \kreuz{quadrierte Klammer – Rückwärtsrechnen} \\ \kreuz{Normalform – Formel}
\end{teile}
\end{aufgabe}

%% TODO Titel: Ich-kann-Satz fehlt in der Bank (Platzhalter: Kettenname) – Nr. 17
%% TODO Anweisung: ein Satz über der ersten Teilaufgabe fehlt in der Bank – Nr. 17
\begin{aufgabe}{Ist die Vorzahl von x² eins?}
\begin{teile}
\teil $4x^2 + 8x - 12 = 0$ – vor der Formel teilen?\\ \kreuz{nicht teilen} \\ \kreuz{durch $4$} \\ \kreuz{durch $8$}
\end{teile}
\end{aufgabe}

%% TODO Titel: Ich-kann-Satz fehlt in der Bank (Platzhalter: Kettenname) – Nr. 18
%% TODO Anweisung: ein Satz über der ersten Teilaufgabe fehlt in der Bank – Nr. 18
\begin{aufgabe}{p-q-Formel}
\begin{teile}
\teil $x^2 - x - 12 = 0$ – p und q? p = \leerfeld , q = \leerfeld
\end{teile}
%% TODO Darstellung für schwach fehlt in der Bank (quadratische-gleichungen-e3-k3-s1-v1; Abschnitt „Für schwache Schüler“)
\swz{$\text{x^2 + 6x + 8 = 0}$}{}
%% TODO Darstellung für schwach fehlt in der Bank (quadratische-gleichungen-e3-k3-s1-v2; Abschnitt „Für schwache Schüler“)
\swz{$\text{x^2 + 10x + 16 = 0}$}{}
%% TODO Darstellung für schwach fehlt in der Bank (quadratische-gleichungen-e3-k3-s1-v3; Abschnitt „Für schwache Schüler“)
\swz{$\text{x^2 + 12x + 20 = 0}$}{}
%% TODO Darstellung für schwach fehlt in der Bank (quadratische-gleichungen-e3-k3-s1-v4; Abschnitt „Für schwache Schüler“)
\swz{$\text{x^2 + 14x + 24 = 0}$}{}
%% TODO Darstellung für schwach fehlt in der Bank (quadratische-gleichungen-e3-k3-s1-v5; Abschnitt „Für schwache Schüler“)
\swz{$\text{x^2 + 10x + 21 = 0}$}{}
\swfrage{Was bleibt gleich, was ändert sich?}
%% TODO Darstellung für schwach fehlt in der Bank (quadratische-gleichungen-e3-k3-s2-v1; Abschnitt „Für schwache Schüler“)
\swz{$\text{x^2 - 4x - 12 = 0}$}{}
%% TODO Darstellung für schwach fehlt in der Bank (quadratische-gleichungen-e3-k3-s3-v1; Abschnitt „Für schwache Schüler“)
\swz{$\text{x^2 + 5 = 6x}$}{}
%% TODO Darstellung für schwach fehlt in der Bank (quadratische-gleichungen-e3-k3-s4-v1; Abschnitt „Für schwache Schüler“)
\swz{$\text{5x^2 - 20x + 15 = 0}$}{}
%% TODO Darstellung für schwach fehlt in der Bank (quadratische-gleichungen-e3-k3-s5-v1; Abschnitt „Für schwache Schüler“)
\swz{$\text{x^2 - 3x + 2 = 0}$}{}
%% TODO Darstellung für schwach fehlt in der Bank (quadratische-gleichungen-e3-k3-s6-v1; Abschnitt „Für schwache Schüler“)
\swz{$\text{x^2 - 10x + 25 = 0}$}{}
%% TODO Darstellung für schwach fehlt in der Bank (quadratische-gleichungen-e3-k3-s7-v1; Abschnitt „Für schwache Schüler“)
\swz{$\text{x^2 + 2x + 5 = 0}$}{}
%% TODO Darstellung für schwach fehlt in der Bank (quadratische-gleichungen-e3-k3-s8-v1; Abschnitt „Für schwache Schüler“)
\swz{$x^2 - 8x + 7 = 0$ – Diskriminante $D = \left(\frac{p}{2}\right)^2 - q$ und Zahl der Lösungen?}{}
%% TODO Darstellung für schwach fehlt in der Bank (quadratische-gleichungen-e3-k3-s9-v1; Abschnitt „Für schwache Schüler“)
\swz{$\text{x^2 - 6x + 4 = 0}$}{}
%% TODO Darstellung für schwach fehlt in der Bank (quadratische-gleichungen-e3-k3-s10-v1; Abschnitt „Für schwache Schüler“)
\swz{Probe: Ist $x = -2$ Lösung von $x^2 - 3x - 10 = 0$?}{}
%% TODO Darstellung für schwach fehlt in der Bank (quadratische-gleichungen-e3-k3-s11-v1; Abschnitt „Für schwache Schüler“)
\swz{$\text{(x + 2)^2 = 5x + 10}$}{}
%% TODO Darstellung für schwach fehlt in der Bank (quadratische-gleichungen-e3-k3-s12-v1; Abschnitt „Für schwache Schüler“)
\swz{$x^2 - 12x = 0$ – welcher Weg, welche Lösungen?}{}
\swz[3]{Die Funktion $f$ mit $f(x) = 2x^2 - 2x - 40$ – Nullstellen? \hfill (P10 2020 OS)}{}
\end{aufgabe}

%% TODO Titel: Ich-kann-Satz fehlt in der Bank (Platzhalter: Kettenname) – Nr. 19
%% TODO Anweisung: ein Satz über der ersten Teilaufgabe fehlt in der Bank – Nr. 19
\begin{aufgabe}{Gerade und Parabel gleichsetzen als Anwendung}
%% TODO Darstellung für schwach fehlt in der Bank (quadratische-gleichungen-e3-k4-s1-v1; Abschnitt „Für schwache Schüler“)
\swz[3]{Parabel $y = x^2$ und Gerade $y = x + 6$ – an welchen Stellen $x$ schneiden sie sich?}{}
\end{aufgabe}

%% TODO Titel: Ich-kann-Satz fehlt in der Bank (Platzhalter: Kettenname) – Nr. 20
%% TODO Anweisung: ein Satz über der ersten Teilaufgabe fehlt in der Bank – Nr. 20
\begin{aufgabe}{p-q-Formel – Fehler finden, Begründen, Darstellungswechsel, Anwendung}
\swz[3]{Wo ist der Fehler? \rechnung{3x^2 + 6x - 24 &= 0 \\ x_{1,2} &= -3 \pm \sqrt{9 + 24}}}{}
\swfrage{Begründe, warum du bei $3x^2 + 9x - 12 = 0$ vor der Formel durch $3$ teilen musst.}
\swz{Die Parabel $y = x^2 - 4x + 3$ ist gezeichnet. Lösungen von $x^2 - 4x + 3 = 0$ am Graphen?}{\begin{ksys}[xmin=-2,xmax=6,ymin=-2,ymax=5,ablesen] \parabel{1}{2}{-1}{} \end{ksys}}
\swz[3]{Ein Ball wird geworfen. Für seine Höhe gilt $h = -0{,}5x^2 + 3x + 8$ ($x$ waagerechte Entfernung in m, $h$ Höhe in m). In welcher Entfernung trifft er auf den Boden?}{}
\end{aufgabe}

%% TODO Merkkasten Einheit 3: 6 Zeilen, Vorgabe höchstens fünf (3.1)
\uebersichtskasten{Normalform: x² + px + q = 0 – alles auf eine Seite, geordnet nach x², x, Zahl; die Vorzahl von x² ist eins, sonst durch sie teilen (normieren). \\ 3x² + 24x − 60 = 0 | : 3 → x² + 8x − 20 = 0 \quad p = 8, q = −20 \\ p-q-Formel: x₁,₂ = −p/2 ± √((p/2)² − q) \\ x² + 8x − 20 = 0 → x₁,₂ = −4 ± √(16 + 20) = −4 ± 6 → x₁ = 2, x₂ = −10 \\ Unter der Wurzel entscheidet: positiv → zwei Lösungen, null → eine Lösung, negativ → keine Lösung. \\ x² + 8x + 16 = 0 → −4 ± √(16 − 16) → x = −4 \quad x² + 8x + 20 = 0 → √(16 − 20): keine Lösung}
